%% 第 17 章：伴随的零空间和值域
%% 本文件由 tools/md2tex.py 自动生成，请勿直接编辑。

\chapter{伴随的零空间和值域}


\begin{jprop}{伴随的零空间}{r:18:1}
\[
\ker T^*
=
(\operatorname{range}T)^\perp.
\]
\end{jprop}

\begin{jproof}{证明思路}
$w\in\ker T^*$，意味着：

\[
T^*(w)=0.
\]

由伴随定义，这等价于：

\[
\langle T(v),w\rangle=0
\]

对所有 $v\in V$ 都成立。

而所有 $T(v)$ 恰好组成 $\operatorname{range}T$。因此，$w$ 与 $T$ 的全部输出正交，即：

\[
w\in(\operatorname{range}T)^\perp.
\]
\end{jproof}

\begin{jprop}{伴随的值域}{r:18:2}
在有限维情形：

\[
\operatorname{range}T^*
=
(\ker T)^\perp.
\]
\end{jprop}

\begin{jproof}{证明思路}
若 $u=T^*(w)$，且 $v\in\ker T$，则：

\[
\langle v,u\rangle
=
\langle v,T^*(w)\rangle
=
\langle T(v),w\rangle
=
0.
\]

因此：

\[
\operatorname{range}T^*
\subseteq
(\ker T)^\perp.
\]

另一方面，伴随与原映射具有相同的秩：

\[
\operatorname{rank}T^*=\operatorname{rank}T.
\]

而：

\[
\dim(\ker T)^\perp
=
\dim V-\dim\ker T
=
\operatorname{rank}T.
\]

一个子空间包含于另一个子空间，且两者维数相同，因此它们相等。
\end{jproof}

\section{四个子空间关系}

伴随将此前的四个子空间关系写成：

\[
\ker T^*
=
(\operatorname{range}T)^\perp,
\]

\[
\operatorname{range}T^*
=
(\ker T)^\perp.
\]

在标准欧氏坐标下，$T^*$ 对应共轭转置；实矩阵情形下就是转置。

这正是经典“四个基本子空间”中行空间、列空间、零空间与左零空间关系的抽象版本，我们将在表示的笔记中更加详细的去说明。
